\newif\ifglobal

%%%%%%%%%%%%%%%%%%%%%%%
%%% Compile to view %%%
%%%%%%%%%%%%%%%%%%%%%%%

%%%%%%%%%%%%%%%%%%%%%%%%%%%%%%%%%%%%%%%%%%%%%%%%%%%
%%% Readme:                                     %%%
%%% https://www.overleaf.com/read/bwdjvfjksvjy/ %%%
%%%%%%%%%%%%%%%%%%%%%%%%%%%%%%%%%%%%%%%%%%%%%%%%%%%

% >>> M?: marks a module setting number or important notice

% >>> M4: Set {./relative-path} to external files for this module <<<
\def \modulefiles {./04-EFUSE}
% To add an external file, use:
% (Replace '---' with actual filename)
% '\input{\modulefiles/tables/---.tex}' for tables
% '\includegraphics{\modulefiles/figures/---}' for figures
% '\subfile{\modulefiles/---}' for register files


% >>> M1: This module is in global project? <<<
% 'YES' - uncomment; 'NO' - comment out
\globaltrue

\ifglobal
    \documentclass[main\_\_CN.tex]{subfiles}

\else
    % Font "FandolSong-Regular" used by ctex does not have GSUB table.
    % It causes warnings that do not affect compilation and can be ignored.
    \PassOptionsToPackage{quiet}{fontspec}

    \documentclass[openany, 10pt]{book}

    % >>> M2: In ./00-shared/preamble-trm-module, also find
    %         settings for visibility of todo notes, watermark <<<
    \usepackage{preamble-trm-module}


    % Enable Chinese typesetting
    \usepackage{ctex}

    % >>> M3: Temporary labels in Chinese
    %         you can add your own labels there <<<
    \myexternaldocument{./00-shared/temp-labels-cn}

\fi %\ifglobal


\setcounter{secnumdepth}{4}


% >>> M5: If this module needs any updates in
% './00-shared/preamble-trm-module.sty'
% (include additional package, etc.), contact your module PM
% or create the file './00-shared/changelog.tex' make the
% necessary updates yourself and document them in this file <<<

% Make text left-aligned and allow latex to break pages naturally, comment out for full justification
\raggedright
\raggedbottom

\begin{document}


% Sets language into which pre-defined bits of text are translated
\selectlanguage{chinese}
\setchineseformatting

% Do not delete this block
\ifglobal
% Add the following front matter if
% this module is outside of global project
\else
    \InputIfFileExists{readme}{}{}
    \listoftodos
    {\let\clearpage\relax \listofdonetodos}
    \tableofcontents
    %\listoftables
    %\listoffigures
\fi



%%%%%%%%%%%%%%%%%%%%%%%%%
%%% TRM Module Begins %%%
%%%%%%%%%%%%%%%%%%%%%%%%%

\hypertarget{efuse}{}
\chapter{eFuse 控制器 (EFUSE)}
\label{mod:efuse}


\section{概述}

ESP32-S2 系统中有一块 4096 位的 eFuse，其中存储着参数内容。eFuse 的各个位 一旦被烧写为 1，则不能再恢复为 0。eFuse 控制器按照用户配置完成对 eFuse 中各参数中的各个位的烧写。从芯片外部，eFuse 数据只能通过 eFuse 控制器读取。对于某些数据，如果未启用读保护，则可以从芯片外部读取该数据；如果启用了读保护，则无法从芯片外部读取该数据。不过，存储在 eFuse 中的某些密钥始终可以供硬件加密模块（例如 RSA\_DS、HMAC 等）在内部使用，芯片外部无法获得这些数据。


\section{主要特性}

\begin{itemize}
\item 一次性可编程存储
    \item 烧写保护可配置
    \item 读取保护可配置
    \item 使用多种硬件编码方式保护参数内容
\end{itemize}


\section{功能描述}

\subsection{结构}

eFuse 从结构上分成 11 个块 (BLOCK0 \verb+~+ BLOCK10)。

BLOCK0 存储大部分核心系统参数，其中 24 位可被用户读取但是对用户没有意义；还有 38 位处于保留状态，留作未来使用。


表 \ref{table:BLOCK0_Parameter} 列出了用户可访问（可读并可用）的所有参数、偏移地址、位宽、是否可供硬件使用、烧写保护，以及功能描述。

在这些参数中，{\bfseries \hyperref[fielddesc:EFUSEWRDIS]{EFUSE\_WR\_DIS}} 用于控制其他参数的烧写，{\bfseries \hyperref[fielddesc:EFUSERDDIS]{EFUSE\_RD\_DIS}} 用于控制用户读取 BLOCK4 \verb+~+ BLOCK10。更多关于这两个参数的信息请见章节 \ref{section:EFUSEWRDISABLE}、\ref{section:EFUSERDDISABLE}。

\begin{longtable}[l]{ |m{5.2cm}|c|c|c|c|m{5.5cm}| }    
\caption{BLOCK0 参数}
    \label{table:BLOCK0_Parameter}
    \\\hline
    \rowcolor{lightgray}
    参数 & \thead{偏移} & \thead{位宽} & \thead{硬件使用} & \thead{\hyperref[fielddesc:EFUSEWRDIS]{EFUSE\_WR\_DIS} \\烧写保护位}
    & 描述 \\
    \endfirsthead
    \hline
    \rowcolor{lightgray}
    参数 & \thead{偏移} & \thead{位宽} & \thead{硬件使用} & \thead{\hyperref[fielddesc:EFUSEWRDIS]{EFUSE\_WR\_DIS} \\烧写保护位}
    & 描述 \\
    \endhead    
    \endfoot
    \endlastfoot
    \hline   
    \hyperref[fielddesc:EFUSEWRDIS]{\textbf{EFUSE\_WR\_DIS}} & 0 & 32 & Y & N/A & 表示是否禁止 eFuse 烧写 \\
    \hline
    \hyperref[fielddesc:EFUSERDDIS]{\textbf{EFUSE\_RD\_DIS}} & 32 & 7 & Y & 0  & 表示是否禁止用户读取 eFuse BLOCK4 \verb+~+ 10 的内容 \\
    \hline
    \hyperref[fielddesc:EFUSEDISICACHE]{EFUSE\_DIS\_ICACHE} & 40 &  1 & Y & 2  & 表示是否关闭 ICache \\
    \hline
    \hyperref[fielddesc:EFUSEDISDCACHE]{EFUSE\_DIS\_DCACHE} & 41 & 1 & Y & 2  & 表示是否关闭 DCache \\
    \hline
    \hyperref[fielddesc:EFUSEDISDOWNLOADICACHE]{EFUSE\_DIS\_DOWNLOAD\_ ICACHE} & 42 & 1 & Y & 2 & 表示是否在 Download 模式下关闭 ICache\\   
    \hline
    \hyperref[fielddesc:EFUSEDISDOWNLOADDCACHE]{EFUSE\_DIS\_DOWNLOAD\_ DCACHE} & 43 & 1 & Y & 2 & 表示是否在 Download 模式下关闭 DCache\\      
    \hline
    \hyperref[fielddesc:EFUSEDISFORCEDOWNLOAD]{EFUSE\_DIS\_FORCE\_DOWNLOAD} & 44 & 1 & Y & 2 & 表示是否禁止强制芯片进入 Download 模式\\   
    \hline
    \hyperref[fielddesc:EFUSEDISUSB]{EFUSE\_DIS\_USB} & 45 &1 & Y & 2 & 表示是否关闭 USB OTG 功能 \\     
    \hline
    \hyperref[fielddesc:EFUSEDISCAN]{EFUSE\_DIS\_CAN} & 46 & 1 & Y & 2 & 表示是否关闭 TWAI 控制器功能 \\  
    \hline
    \hyperref[fielddesc:EFUSEDISBOOTREMAP]{EFUSE\_DIS\_BOOT\_REMAP} & 47 & 1 & Y & 2 & 表示 RAM 空间是否可以映射到 ROM 空间 \\   
    \hline
    \hyperref[fielddesc:EFUSESOFTDISJTAG]{EFUSE\_SOFT\_DIS\_JTAG} & 49 & 1 & Y & 2 & 表示是否软禁用 JTAG \\    
    \hline
    \hyperref[fielddesc:EFUSEHARDDISJTAG]{EFUSE\_HARD\_DIS\_JTAG} & 50 & 1 & Y & 2 & 表示是否硬禁用 JTAG（永久） \\   
    \hline
    \hyperref[fielddesc:EFUSEDISDOWNLOADMANUALENCRYPT]{EFUSE\_DIS\_DOWNLOAD\_ MANUAL\_ENCRYPT} & 51 &  1 & Y & 2 & 表示是否在 download boot 模式下禁用 flash 加密功能 \\ 
    \hline
    \hyperref[fielddesc:EFUSEUSBEXCHGPINS]{EFUSE\_USB\_EXCHG\_PINS} & 56 & 1 & Y & 30 & 表示是否交换 USB D+/D- 管脚 \\
    \hline
    \hyperref[fielddesc:EFUSEEXTPHYENABLE]{EFUSE\_EXT\_PHY\_ENABLE} & 57 & 1 & N & 30 & 表示是否使能外部 USB PHY \\
    \hline
    \hyperref[fielddesc:EFUSEUSBFORCENOPERSIST]{EFUSE\_USB\_FORCE\_NOPERSIST} & 58 & 1 & N & 30 & 表示是否强制将 USB BVALID 设置为 1 \\    
    \hline
    \hyperref[fielddesc:EFUSEVDDSPIXPD]{EFUSE\_VDD\_SPI\_XPD} & 68 & 1 & Y & 3 & 表示 VDD\_SPI\_FORCE 为 1 时，VDD\_SPI 调节器是否上电 \\        
    \hline
    \hyperref[fielddesc:EFUSEVDDSPITIEH]{EFUSE\_VDD\_SPI\_TIEH} & 69 & 1 & Y & 3 & 表示 VDD\_SPI\_FORCE 为 1 时，VDD\_SPI 的电压选择 \\   
    \hline
    \hyperref[fielddesc:EFUSEVDDSPIFORCE]{EFUSE\_VDD\_SPI\_FORCE} & 70 & 1 & Y & 3 & 表示是否使用 XPD\_VDD\_PSI\_REG 和 VDD\_SPI\_TIEH 配置 VDD\_SPI LDO \\ 
    \hline
    \hyperref[fielddesc:EFUSEWDTDELAYSEL]{EFUSE\_WDT\_DELAY\_SEL} & 80 & 2 & Y & 3 & 表示是否选择 RTC 看门狗超时阈值 \\
    \hline
    \hyperref[fielddesc:EFUSESPIBOOTCRYPTCNT]{EFUSE\_SPI\_BOOT\_CRYPT\_CNT} & 82 & 3 & Y & 4 & 表示是否使能 SPI boot 加解密 \\
    \hline
    \hyperref[fielddesc:EFUSESECUREBOOTKEYREVOKE0]{EFUSE\_SECURE\_BOOT\_KEY\_ REVOKE0} & 85 & 1 & N & 5 & 表示是否撤销第一个安全启动 (Secure Boot) 密钥 \\
    \hline
    \hyperref[fielddesc:EFUSESECUREBOOTKEYREVOKE1]{EFUSE\_SECURE\_BOOT\_KEY\_ REVOKE1} & 86 & 1 & N & 6 & 表示是否撤销第二个安全启动密钥 \\
    \hline
    \hyperref[fielddesc:EFUSESECUREBOOTKEYREVOKE2]{EFUSE\_SECURE\_BOOT\_KEY\_ REVOKE2} & 87 & 1 & N & 7 & 表示是否撤销第三个安全启动密钥 \\
    \hline
    \hyperref[fielddesc:EFUSEKEYPURPOSE0]{EFUSE\_KEY\_PURPOSE\_0} & 88 & 4 & Y & 8 & 表示 Key0 用途 (purpose)，见表 \ref{table: key_purpose} \\
    \hline
    \hyperref[fielddesc:EFUSEKEYPURPOSE1]{EFUSE\_KEY\_PURPOSE\_1} & 92 & 4 & Y & 9 & 表示 Key1 用途，见表 \ref{table: key_purpose} \\
    \hline
    \hyperref[fielddesc:EFUSEKEYPURPOSE2]{EFUSE\_KEY\_PURPOSE\_2} & 96 & 4 & Y & 10 & 表示 Key2 用途，见表 \ref{table: key_purpose} \\ 
    \hline
    \hyperref[fielddesc:EFUSEKEYPURPOSE3]{EFUSE\_KEY\_PURPOSE\_3} & 100 & 4 & Y & 11 &  表示 Key3 用途，见表 \ref{table: key_purpose} \\    
    \hline
    \hyperref[fielddesc:EFUSEKEYPURPOSE4]{EFUSE\_KEY\_PURPOSE\_4} & 104 & 4 & Y & 12 & 表示 Key4 用途，见表 \ref{table: key_purpose} \\    
    \hline
    \hyperref[fielddesc:EFUSEKEYPURPOSE5]{EFUSE\_KEY\_PURPOSE\_5} & 108 & 4 & Y & 13 & 表示 Key5 用途，见表 \ref{table: key_purpose} \\  
    \hline
    \hyperref[fielddesc:EFUSESECUREBOOTEN]{EFUSE\_SECURE\_BOOT\_EN} & 116 & 1 & N & 15 & 表示是否使能安全启动 \\ 
    \hline
    \hyperref[fielddesc:EFUSESECUREBOOTAGGRESSIVEREVOKE]{EFUSE\_SECURE\_BOOT\_ AGGRESSIVE\_REVOKE} & 117 & 1 & N & 16 & 表示是否使安全启动的撤销采用激进策略 \\
    \hline
    \hyperref[fielddesc:EFUSEFLASHTPUW]{EFUSE\_FLASH\_TPUW} & 124 & 4 & N & 18 & 表示上电后 flash 等待时间 \\    
    \hline    
    \hyperref[fielddesc:EFUSEDISDOWNLOADMODE]{EFUSE\_DIS\_DOWNLOAD\_MODE} & 128 & 1 & N & 18 & 表示是否关闭所有 download boot 模式\\       
    \hline
    \hyperref[fielddesc:EFUSEDISLEGACYSPIBOOT]{EFUSE\_DIS\_LEGACY\_SPI\_BOOT} & 129 & 1 & N & 18 & 表示是否关闭 Legacy SPI boot 模式 \\
    \hline
    \hyperref[fielddesc:EFUSEUARTPRINTCHANNEL]{EFUSE\_UART\_PRINT\_CHANNEL} & 130 & 1 & N & 18 & 表示打印 boot 信息的 UART 通道选择 \\
    \hline 
    \hyperref[fielddesc:EFUSEDISUSBDOWNLOADMODE]{EFUSE\_DIS\_USB\_DOWNLOAD\_ MODE} & 132 & 1 & N & 18 & 表示是否在 UART download boot 模式下关闭 USB OTG 下载功能 \\
    \hline 
    \hyperref[fielddesc:EFUSEENABLESECURITYDOWNLOAD]{EFUSE\_ENABLE\_SECURITY\_ DOWNLOAD} & 133 & 1 & N & 18 & 表示是否使能 UART 安全下载模式 \\
    \hline
    \hyperref[fielddesc:EFUSEUARTPRINTCONTROL]{EFUSE\_UART\_PRINT\_CONTROL} & 134 & 2 & N & 18 & 表示 UART boot 信息的打印方式 \\
    \hline
    \hyperref[fielddesc:EFUSEPINPOWERSELECTION]{EFUSE\_PIN\_POWER\_SELECTION} & 136 & 1 & N & 18 & 表示 GPIO33 \verb+~+ GPIO37 的电源选择 \\
    \hline
    \hyperref[fielddesc:EFUSEFLASHTYPE]{EFUSE\_FLASH\_TYPE} & 137 & 1 & N & 18 & 表示 flash 类型 \\
    \hline
    \hyperref[fielddesc:EFUSEFORCESENDRESUME]{EFUSE\_FORCE\_SEND\_RESUME} & 138 & 1 & N & 18 & 表示是否强制 ROM 代码在 SPI 启动过程中发送恢复指令 \\  
    \hline
    \hyperref[fielddesc:EFUSESECUREVERSION]{EFUSE\_SECURE\_VERSION} & 139 & 16 & N & 18 & 表示 IDF 安全版本 \\\hline
\end{longtable}


表 \ref{table: key_purpose} 为密钥用途各个数值对应的含义。通过配置参数 EFUSE\_KEY\_PURPOSE\_\textcolor{red}{\it{n}} 来声明 KEY\textcolor{red}{\it{n}} 用途 (\textcolor{red}{\it{n}}: 0 \verb+~+ 5)。

\begin{longtable}[c]{|c|p{10cm}|}
\caption{密钥用途数值对应的含义}
\label{table: key_purpose}
    \\\hline
    \rowcolor{lightgray}
  密钥用途数值  & 含义\\
    \endfirsthead
    \hline
    \rowcolor{lightgray}
   密钥用途数值  & 含义\\
    \endhead
    \endfoot
    \endlastfoot
    \hline
0 & 指定为用户使用 \\\hline
1&  保留 \\\hline
2& 指定为 XTS\_AES\_256\_KEY\_1 使用（用于 flash/SRAM 加解密） \\\hline
3& 指定为 XTS\_AES\_256\_KEY\_2 使用（用于 flash/SRAM 加解密） \\ \hline
4 & 指定为 XTS\_AES\_128\_KEY 使用（用于 flash/SRAM 加解密） \\\hline
5& 指定为 HMAC Downstream（下行）模式使用\\\hline
6& 指定为 HMAC Downstream 模式下的 JTAG 使用 \\\hline
7 & 指定为 HMAC Downstream 模式下的 RSA\_DS 使用\\\hline
8 &指定为 HMAC Upstream（上行）模式使用\\\hline
9& 指定为 SECURE\_BOOT\_DIGEST0 使用（secure boot 密钥摘要） \\\hline
10 &指定为 SECURE\_BOOT\_DIGEST1 使用（secure boot 密钥摘要）\\\hline
11& 指定为 SECURE\_BOOT\_DIGEST2 使用（secure boot 密钥摘要）\\\hline
\end{longtable}

表 \ref{table:BLOCK1_10_Parameter} 列出了 BLOCK1 \verb+~+ BLOCK10 中存储的参数的信息。

\begin{longtable}[l]{ |c|c|c|c|c|c|l| }
    \caption{BLOCK1-10 参数}
    \label{table:BLOCK1_10_Parameter}
    \\\hline
    \rowcolor{lightgray}
    块 & 参数 & 位宽 & 硬件使用 &  \thead{\hyperref[fielddesc:EFUSEWRDIS]{EFUSE\_WR\_DIS} \\ 烧写保护位} &
   \thead{ \hyperref[fielddesc:EFUSEWRDIS]{EFUSE\_RD\_DIS} \\ 用户读取保护位} &
    描述 \\
    \endfirsthead
    \hline
    \rowcolor{lightgray}
    块 & 参数 & 位宽 & 硬件使用 & \thead{\hyperref[fielddesc:EFUSERDDIS]{EFUSE\_WR\_DIS} \\ 烧写保护位} &
   \thead{ \hyperref[fielddesc:EFUSERDDIS]{EFUSE\_RD\_DIS} \\ 用户读取保护位} &
    描述 \\\hline
    \endhead
    \endfoot
    \endlastfoot
    \hline
    BLOCK1 & \hyperref[fielddesc:EFUSEMAC0]{EFUSE\_MAC} & 48 & N & 20 & N/A & MAC 地址 \\\cline{2-7}
    & \hyperref[fielddesc:EFUSESPIPADCONF0]{EFUSE\_SPI\_PAD\_}& [0:5] & N & 20 & N/A & {CLK} \\
    \cline{3-7}
    & \hyperref[fielddesc:EFUSESPIPADCONF0]{CONFIGURE} & [6:11] & N & 20 & N/A & Q (D1) \\
    \cline{3-7}
    & & [12:17] & N & 20 & N/A & D (D0) \\
    \cline{3-7}
    & & [18:23] & N & 20 & N/A & CS\\
    \cline{3-7}
    & & [24:29] & N & 20 & N/A & HD (D3) \\
    \cline{3-7}
    & & [30:35] & N & 20 & N/A & WP (D2) \\
    \cline{3-7}
    & & [36:41] & N & 20 & N/A & DQS \\
    \cline{3-7}
    & & [42:47] & N & 20 & N/A & D4 \\
    \cline{3-7}
    & & [48:53] & N & 20 & N/A & D5 \\
    \cline{3-7}
    & & [54:59] & N & 20 & N/A & D6 \\
    \cline{3-7}
    & & [60:65] & N & 20 & N/A & D7 \\
    \cline{ 2-7 }
    & \hyperref[fielddesc:EFUSESYSDATAPART00]{EFUSE\_SYS\_DATA\_PART0} & 78 & N & 20 & N/A & 系统数据 \\\hline		
    BLOCK2 & \hyperref[fielddesc:EFUSESYSDATAPART1N]{EFUSE\_SYS\_DATA\_PART1} & 256 & N & 21 & N/A & 系统数据 \\\hline		
    BLOCK3 & \hyperref[fielddesc:EFUSEUSRDATAN]{EFUSE\_USR\_DATA} & 256 & N & 22 & N/A & 用户数据 \\\hline
    BLOCK4 & \hyperref[fielddesc:EFUSEKEY0DATAN]{EFUSE\_KEY0\_DATA} & 256 & Y & 23 & 0 & KEY0 或用户数据 \\
    \hline
    BLOCK5 & \hyperref[fielddesc:EFUSEKEY1DATAN]{EFUSE\_KEY1\_DATA} & 256 & Y & 24 & 1 & KEY1 或用户数据 \\
    \hline
    BLOCK6 & \hyperref[fielddesc:EFUSEKEY2DATAN]{EFUSE\_KEY2\_DATA} & 256 & Y & 25 & 2 & KEY2 或用户数据 \\    
    \hline
    BLOCK7 & \hyperref[fielddesc:EFUSEKEY3DATAN]{EFUSE\_KEY3\_DATA} & 256 & Y & 26 & 3 & KEY3 或用户数据 \\      
    \hline
    BLOCK8 & \hyperref[fielddesc:EFUSEKEY4DATAN]{EFUSE\_KEY4\_DATA} & 256 & Y & 27 & 4 & KEY4 或用户数据 \\      
    \hline
    BLOCK9 & \hyperref[fielddesc:EFUSEKEY5DATAN]{EFUSE\_KEY5\_DATA} & 256 & Y & 28 & 5 & KEY5 或用户数据 \\      
    \hline
    BLOCK10 & \hyperref[fielddesc:EFUSESYSDATAPART2N]{EFUSE\_SYS\_DATA\_PART2} & 256 & N & 29 & 6 & 系统数据 \\ 
    \hline
\end{longtable}

其中，BLOCK4 \verb+~+ 9 分别存储 KEY0\verb+~+5，表示 eFuse 中至多可以烧写 6 个 256 位的密钥。每烧写一个密钥，还需要烧写该密钥用途的数值（见表 \ref{table: key_purpose}）。例如，用户将用于 HMAC Downstream 模式下的 JTAG 功能的密钥烧写到 KEY\textcolor{red}{3}（即 BLOCK7），还需要将密钥用途的数值 6 烧写到 \hyperref[fielddesc:EFUSEKEYPURPOSE3]{EFUSE\_KEY\_PURPOSE\_\textcolor{red}{3}}。

BLOCK1 \verb+~+ BLOCK10 均采用 RS 编码方式，因此参数烧写受到一定的限制，具体请参考章节 \ref{datastorage}：\textit{数据存储方式}，和章节 \ref{Software Programming of Parameters}：\textit{烧写参数}。

\subsubsection{EFUSE\_WR\_DIS}\label{section:EFUSEWRDISABLE}

参数 \hyperref[fielddesc:EFUSEWRDIS]{EFUSE\_WR\_DIS} 决定了 eFuse 中所有的参数是否处于烧写保护状态。烧写完 \hyperref[fielddesc:EFUSEWRDIS]{EFUSE\_WR\_DIS} 参数后，需要更新 eFuse 读寄存器才能生效（参考章节 \ref{updatingRDregister} 中：\textit{更新 eFuse 读寄存器}）。

表 \ref{table:BLOCK0_Parameter} 以及表 \ref{table:BLOCK1_10_Parameter} 中的“\hyperref[fielddesc:EFUSEWRDIS]{EFUSE\_WR\_DIS} 烧写保护位”列描述了各参数的烧写保护状态具体由 \hyperref[fielddesc:EFUSEWRDIS]{EFUSE\_WR\_DIS} 的哪个位决定。

当某个参数对应的烧写保护位为 0 时，表示此参数未处于烧写保护状态，可以烧写该参数，但已经被烧写的参数不能被重复烧写。

当某个参数对应的烧写保护位为 1 时，表示此参数处于烧写保护状态，此参数的每一个位都无法被更改，未被烧写的位永远为 0，已经被烧写的位永远为 1。所以如果某个参数已经处于烧写保护状态了，则会一直处在该状态，无法再更改。

\subsubsection{EFUSE\_RD\_DIS}\label{section:EFUSERDDISABLE}

所有参数中，只有 BLOCK4 \verb+~+ BLOCK10 的参数受用户读取保护状态的约束，即表 \ref{table:BLOCK1_10_Parameter} 中“\hyperref[fielddesc:EFUSERDDIS]{EFUSE\_RD\_DIS} 用户读取保护”列非“N/A”的参数。烧写完 \hyperref[fielddesc:EFUSERDDIS]{EFUSE\_RD\_DIS} 参数后，需要更新 eFuse 读寄存器才能生效（参考章节 \ref{updatingRDregister} 中：\textit{更新 eFuse 读寄存器}）。

参数 \hyperref[fielddesc:EFUSERDDIS]{EFUSE\_RD\_DIS} 中的某个位为 0，表示此位管理的参数未处于用户读取保护状态；某个位为 1，表示此位管理的参数处于用户读取保护状态。

除 BLOCK4 \verb+~+ BLOCK10 之外，其他参数不受用户读取保护状态的约束，均可被用户读取。

BLOCK4 \verb+~+ BLOCK10 即使被配置处于读取保护状态，仍然可以通过设置 \hyperref[fielddesc:EFUSEKEYPURPOSE3]{EFUSE\_KEY\_PURPOSE\_\textcolor{red}{\it{n}}} 供硬件加密模块在内部使用。


\subsubsection{数据存储方式}\label{datastorage}

eFuse 使用硬件编码机制保护数据，对用户不可见。

BLOCK0 使用 4 备份方式存储参数，即 BLOCK0 中的所有参数（除了 \hyperref[fielddesc:EFUSEWRDIS]{EFUSE\_WR\_DIS} ）均在 eFuse 中存储了 4 份。4 备份机制对用户不可见。

BLOCK1 \verb+~+ BLOCK10 使用 RS (44, 32) 编码方式，最多支持自动校正 6 个字节。本文 RS (44, 32) 使用的本源多项式为 \(p(x)=x^8+x^4+x^3+x^2+1\)。

\begin{figure}[H]
    \centering
    \fbox{\includegraphics[width=1.0\textwidth]{\modulefiles/figures/lfsr_circuit_20190809_1}}
    \caption{移位寄存器电路图（前 32 字节）}
    \label{Figure:lfsr_circuit_1}
\end{figure}

\begin{figure}[H]
    \centering
    \fbox{\includegraphics[width=1.0\textwidth]{\modulefiles/figures/lfsr_circuit_20190809_2}}
    \caption{移位寄存器电路图（后 12 字节）}
    \label{Figure:lfsr_circuit_2}
\end{figure}

如图 \ref{Figure:lfsr_circuit_1} 和 \ref{Figure:lfsr_circuit_2}所示，用户对 32 字节参数进行 RS (44, 32) 编码处理，将 32 字节数据处理为 44 字节，其中：

\begin{itemize}
    \item 字节 [0:31] 为数据本身
    \item 字节 [32:43] 为储存在 8 位触发器 DFF1, DFF2, ..., DFF12 中的奇偶校验字节（gf\_mul\_n 为 \(GF(2^8)\) 域中某一字节数据与元素 \(\alpha^n\) 相乘的结果，n 为整数）
\end{itemize}

然后，硬件将这 44 字节数据一起烧入 eFuse。eFuse 控制器会在读 eFuse 的过程中自动完成解码和自动校正。

由于 RS 校验码是在整个 256 位的 eFuse block 上生成的，因此每个 block 只能写入一次。

\subsection{烧写参数}\label{Software Programming of Parameters}


烧写 eFuse 参数时，需要按块烧写。BLOCK0 \verb+~+ BLOCK10 共用同一段地址来存储即将烧写的参数。通过配置 \hyperref[fielddesc:EFUSEBLKNUM]{EFUSE\_BLK\_NUM} 参数表明当前需要烧写的是哪一个块。

{\bfseries 烧写 BLOCK0}

当 \hyperref[fielddesc:EFUSEBLKNUM]{EFUSE\_BLK\_NUM} = 0 时，烧写 BLOCK0。\hyperref[regdesc:EFUSEPGMDATANREG]{EFUSE\_PGM\_DATA0\_REG} 寄存器存储着 \hyperref[fielddesc:EFUSEWRDIS]{EFUSE\_WR\_DIS}。\hyperref[regdesc:EFUSEPGMDATANREG]{EFUSE\_PGM\_DATA1\_REG} \verb+~+ \hyperref[regdesc:EFUSEPGMDATANREG]{EFUSE\_PGM\_DATA5\_REG} 用来存储即将烧写的参数的有效信息，其中 24 位为用户可读但对用户没有意义的有效信息，必须写入 0，对应位置为：

\begin{itemize}
    \item \hyperref[regdesc:EFUSEPGMDATANREG]{EFUSE\_PGM\_DATA1\_REG}[29:31]
    \item \hyperref[regdesc:EFUSEPGMDATANREG]{EFUSE\_PGM\_DATA1\_REG}[20:23]
    \item \hyperref[regdesc:EFUSEPGMDATANREG]{EFUSE\_PGM\_DATA2\_REG}[7:15]
    \item \hyperref[regdesc:EFUSEPGMDATANREG]{EFUSE\_PGM\_DATA2\_REG}[0:3]
    \item \hyperref[regdesc:EFUSEPGMDATANREG]{EFUSE\_PGM\_DATA3\_REG}[16:19]
\end{itemize}


\hyperref[regdesc:EFUSEPGMDATANREG]{EFUSE\_PGM\_DATA6\_REG} \verb+~+ \hyperref[regdesc:EFUSEPGMDATANREG]{EFUSE\_PGM\_DATA7\_REG} 以及 \hyperref[regdesc:EFUSEPGMCHECKVALUENREG]{EFUSE\_PGM\_CHECK\_VALUE0\_REG} \verb+~+ \hyperref[regdesc:EFUSEPGMCHECKVALUENREG]{EFUSE\_PGM\_CHECK\_VALUE2\_REG} 中的数据不影响 BLOCK0 的烧写。

{\bfseries 烧写 BLOCK1}

当 \hyperref[fielddesc:EFUSEBLKNUM]{EFUSE\_BLK\_NUM} = 1 时，\hyperref[regdesc:EFUSEPGMDATANREG]{EFUSE\_PGM\_DATA0\_REG} \verb+~+ \hyperref[regdesc:EFUSEPGMDATANREG]{EFUSE\_PGM\_DATA5\_REG} 存储着 BLOCK1 即将烧写的参数，\hyperref[regdesc:EFUSEPGMCHECKVALUENREG]{EFUSE\_PGM\_CHECK\_VALUE0\_REG} \verb+~+ \hyperref[regdesc:EFUSEPGMCHECKVALUENREG]{EFUSE\_PGM\_CHECK\_VALUE2\_REG} 中存储着对应的 RS 校验码。\hyperref[regdesc:EFUSEPGMDATANREG]{EFUSE\_PGM\_DATA6\_REG} \verb+~+ \hyperref[regdesc:EFUSEPGMDATANREG]{EFUSE\_PGM\_DATA7\_REG} 中的数据不影响 BLOCK1 的烧写。RS 校验码的计算视这些位为 0。

{\bfseries 烧写 BLOCK2 \verb+~+ 10}

当 \hyperref[fielddesc:EFUSEBLKNUM]{EFUSE\_BLK\_NUM} = 2 \verb+~+ 10 时，\hyperref[regdesc:EFUSEPGMDATANREG]{EFUSE\_PGM\_DATA0\_REG} \verb+~+ \hyperref[regdesc:EFUSEPGMDATANREG]{EFUSE\_PGM\_DATA7\_REG} 存储着即将烧写的参数，\hyperref[regdesc:EFUSEPGMCHECKVALUENREG]{EFUSE\_PGM\_CHECK\_VALUE0\_REG}  \verb+~+ \hyperref[regdesc:EFUSEPGMCHECKVALUENREG]{EFUSE\_PGM\_CHECK\_VALUE2\_REG} 中存储着对应的 RS 校验码。

{\bfseries 烧写流程}

烧写参数的流程如下：
\begin{enumerate}
\item 配置 \hyperref[fielddesc:EFUSEBLKNUM]{EFUSE\_BLK\_NUM} 参数，决定烧写哪一个块。

\item 将需要烧写的参数填写到寄存器 \hyperref[regdesc:EFUSEPGMDATANREG]{EFUSE\_PGM\_DATA0\_REG} \verb+~+ \hyperref[regdesc:EFUSEPGMDATANREG]{EFUSE\_PGM\_DATA7\_REG} 和 \hyperref[regdesc:EFUSEPGMCHECKVALUENREG]{EFUSE\_PGM\_CHECK\_VALUE0\_REG}  \verb+~+  \hyperref[regdesc:EFUSEPGMCHECKVALUENREG]{EFUSE\_PGM\_CHECK\_VALUE2\_REG} 中。

\item 确保 eFuse 时钟寄存器的配置正确，具体请参考章节 \ref{eFuse-Programming Timing}：\textit{eFuse
烧写时序}。

\item 确保 eFuse 烧写电压 VDDQ 的配置正确，具体请参考章节 \ref{eFuse VDDQ Timing}：\textit{eFuse VDDQ
时序}。

\item 配置寄存器 \hyperref[regdesc:EFUSECONFREG]{EFUSE\_CONF\_REG} 的 \hyperref[fielddesc:EFUSEOPCODE]{EFUSE\_OP\_CODE} 位段为 \hexprefix{}5A5A。

\item 配置寄存器 \hyperref[regdesc:EFUSECMDREG]{EFUSE\_CMD\_REG} 的 \hyperref[fielddesc:EFUSEPGMCMD]{EFUSE\_PGM\_CMD} 位段为 1。

\item 轮询寄存器 \hyperref[regdesc:EFUSECMDREG]{EFUSE\_CMD\_REG} 直到其为 \hexprefix{}0，或者等待烧写完成中断产生。识别烧写/读取完成中断产生的方法详见章节 \ref{softwarereadparameters} 最后的说明。

\item 将寄存器中写入的参数清零。

\item 执行更新 eFuse 读寄存器操作使写入的新值生效，具体请参考章节 \ref{softwarereadparameters}：\textit{用户读取参数}。
\end{enumerate}

{\bfseries 限制}

BLOCK0 中不同的参数，甚至对于同一个参数中的不同位可以在多次烧写中分别完成。
但是并不推荐这样做，而是建议尽量减少烧写次数。
我们建议对于某个参数中的所有需要烧写的位都在一次烧写中完成。
并且当 \hyperref[fielddesc:EFUSEWRDIS]{EFUSE\_WR\_DIS} 的某个位管理的所有参数都烧写之后，
就立即烧写 \hyperref[fielddesc:EFUSEWRDIS]{EFUSE\_WR\_DIS} 的这个位。
甚至可以在同一次烧写中既烧写 \hyperref[fielddesc:EFUSEWRDIS]{EFUSE\_WR\_DIS} 的某个位管理的所有参数，
同时也烧写 \hyperref[fielddesc:EFUSEWRDIS]{EFUSE\_WR\_DIS} 的这个位。另外严禁对已经烧写了的位重复烧写，否则将发生烧写错误。

BLOCK1 中数据信息在出厂时已经烧写完毕，不允许再次烧写。

BLOCK2 \verb+~+ 10 中每一个 BLOCK 都只能烧写一次，不允许重复烧写。


\subsection{用户读取参数}\label{softwarereadparameters}

用户不能直接读取 eFuse 中烧写的信息内容。eFuse 控制器能够将烧写的数据信息读取到对应的地址段的寄存器内，用户再通过读取以 EFUSE\_RD\_ 开始的寄存器来获取 eFuse 信息。表 \ref{fig:efuse-read-register} 列出了读取数据的寄存器名称以及对应烧写时的烧写寄存器名称。


\begin{longtable}[c]{|c|l|l|}
    \caption{用户读取参数相关寄存器} \label{fig:efuse-read-register}
    \\\hline
    \rowcolor{lightgray}
  BLOCK & 读寄存器 & 烧写寄存器\\
    \endfirsthead
    \hline
    \rowcolor{lightgray}
  BLOCK & 读寄存器 & 烧写寄存器\\
    \endhead
    \endfoot
    \endlastfoot
    \hline
0 & \hyperref[regdesc:EFUSERDWRDISREG]{EFUSE\_RD\_WR\_DIS\_REG} & \hyperref[regdesc:EFUSEPGMDATANREG]{EFUSE\_PGM\_DATA0\_REG} \\\hline
0 & EFUSE\_RD\_REPEAT\_DATA0 \verb+~+ 4\_REG & EFUSE\_PGM\_DATA1 \verb+~+ 5\_REG \\\hline
1 & EFUSE\_RD\_MAC\_SPI\_SYS\_0 \verb+~+ 5\_REG & EFUSE\_PGM\_DATA0 \verb+~+ 5\_REG \\\hline
2 & EFUSE\_RD\_SYS\_PART1\_0 \verb+~+ 7\_REG & EFUSE\_PGM\_DATA0 \verb+~+ 7\_REG \\\hline
3 & EFUSE\_RD\_USR\_DATA0 \verb+~+ 7\_REG & EFUSE\_PGM\_DATA0 \verb+~+ 7\_REG \\ \hline
4-9 & EFUSE\_RD\_KEY\textcolor{red}{\it{n}}\_DATA0 \verb+~+ 7\_REG (\textcolor{red}{\it{n}}: 0 \verb+~+ 5) & EFUSE\_PGM\_DATA0 \verb+~+ 7\_REG \\\hline
10 & EFUSE\_RD\_SYS\_PART2\_0 \verb+~+ 7\_REG & EFUSE\_PGM\_DATA0 \verb+~+ 7\_REG \\\hline
\end{longtable}

{\bfseries 更新 eFuse 读寄存器}\label{updatingRDregister}

eFuse 控制器读取内部 eFuse 来更新相应寄存器的数据。读取操作在系统复位时进行，也可以根据需要由用户手动触发（例如在需要读取新烧写 eFuse 中的数据内容时）。用户触发 eFuse 读取操作的流程如下：
\begin{enumerate}
    \item 配置 eFuse 读取时序，具体请参考章节 \ref{eFuse-Read Timing}：\textit{eFuse
读取时序}。
    \item 配置寄存器 \hyperref[regdesc:EFUSECONFREG]{EFUSE\_CONF\_REG} 的 \hyperref[fielddesc:EFUSEOPCODE]{EFUSE\_OP\_CODE} 位段为 \hexprefix{}5AA5。
    \item 配置寄存器 \hyperref[regdesc:EFUSECMDREG]{EFUSE\_CMD\_REG} 的 \hyperref[fielddesc:EFUSEREADCMD]{EFUSE\_READ\_CMD} 位段为 1。
    \item 轮询寄存器 \hyperref[regdesc:EFUSECMDREG]{EFUSE\_CMD\_REG} 直到其为 \hexprefix{}0，或者等待 read\_done interrupt（读取完成中断）产生，识别烧写/读取完成中断产生的方法详见下方说明。
    \item 用户从 eFuse 存储器中读取参数的值。
\end{enumerate}
eFuse 读寄存器中的数值将一直保持到下一次执行更新 eFuse 读操作。

{\bfseries 烧写错误检测}

烧写错误记录寄存器允许用户检测 eFuse 参数的备份是否有不一致的错误。

EFUSE\_RD\_REPEAT\_ERR0 \verb+~+ 3\_REG 寄存器用于指示 BLOCK0 中除了 \hyperref[fielddesc:EFUSEWRDIS]{EFUSE\_WR\_DIS} 外的其他参数的烧写是否出错（对应位为 1 代表烧写出错，此位作废；为 0 代表烧写正确）。

EFUSE\_RD\_RS\_ERR0 \verb+~+ 1\_REG 寄存器记录 eFuse 读 BLOCK1 \verb+~+ BLOCK10 过程中，纠错的字节数目以及 RS 解码是否失败的信息。

用户只能在更新 eFuse 读寄存器操作完成之后才可以读取上面几个寄存器的值。

{\bfseries 识别烧写/读取操作完成}\label{tip:program and read done interrupt}

识别烧写/读取操作完成的方法如下。位 1 对应烧写操作，位 0 对应读取操作。
\begin{itemize}

    \item 方法 1：
        \begin{enumerate}
            \item 轮询寄存器 \hyperref[regdesc:EFUSEINTRAWREG]{EFUSE\_INT\_RAW\_REG} 的位 1/0，直到位 1/0 为 1，表示烧写/读取操作完成。
        \end{enumerate}
    \item 方法2:
        \begin{enumerate}
            \item 对寄存器 \hyperref[regdesc:EFUSEINTENAREG]{EFUSE\_INT\_ENA\_REG} 的位 1/0 置 1，使 eFuse 控制器能够产生烧写/读取完成中断。
            \item 配置中断矩阵使 CPU 能够响应 EFUSE 的中断信号。

            \item 等待烧写/读取完成中断产生。

            \item 对寄存器 \hyperref[regdesc:EFUSEINTCLRREG]{EFUSE\_INT\_CLR\_REG} 的位 1/0 置 1 以清除烧写/读取完成中断。
    \end{enumerate}
\end{itemize}

\subsection{时序}
\subsubsection{eFuse 烧写时序}\label{eFuse-Programming Timing}
图 \ref{Figure:efuse_pgm} 是 eFuse 烧写的时序图。硬件提供 \hyperref[fielddesc:EFUSETSUPA]{EFUSE\_TSUP\_A}、\hyperref[fielddesc:EFUSETPGM]{EFUSE\_TPGM}、\hyperref[fielddesc:EFUSETHPA]{EFUSE\_THP\_A} 和 \hyperref[fielddesc:EFUSETPGMINACTIVE]{EFUSE\_TPGM\_INACTIVE} 四个寄存器用于控制 eFuse 的烧写时序。图中 CSB、VDDQ、PGENB 的含义如下：
\begin{itemize}
\item CSB：片选信号，低电平有效
\item VDDQ：eFuse烧写电压
\item PGENB：烧写使能信号，低电平有效
\end{itemize}

\begin{figure}[H]
    \centering
    \fbox{\includegraphics[width=0.9\textwidth]{\modulefiles/figures/efuse_pgm_20191025}}
    \caption{eFuse 烧写时序图}
    \label{Figure:efuse_pgm}
\end{figure}

eFuse 模块使用的是 CLK\_APB 时钟。由于 CLK\_APB 时钟频率是可变的，需要根据不同的时钟频率来配置上述时序参数，具体请见表 \ref{table:Burn_Timing_Configuration}。复位后，eFuse 默认的参数配置对应的是 20 MHz 时钟频率。

\begin{longtable}[c]{ |c|c|c|c|c|}
    \caption{eFuse 烧写时序参数配置}
    \label{table:Burn_Timing_Configuration}
    \\\hline
    \rowcolor{lightgray}
    \thead{APB Frequency} & \thead{\hyperref[fielddesc:EFUSETSUPA]{EFUSE\_TSUP\_A} \\ (> 6.669 ns)} & \thead{\hyperref[fielddesc:EFUSETPGM]{EFUSE\_TPGM} \\ (9-11 $\mu$s, usually 10 $\mu$s)} & \thead{\hyperref[fielddesc:EFUSETHPA]{EFUSE\_THP\_A} \\ (> 6.166 ns)} & \thead{\hyperref[fielddesc:EFUSETPGMINACTIVE]{EFUSE\_TPGM\_INACTIVE} \\ (> 35.96 ns)} \\
    \endfirsthead
    \hline
    \rowcolor{lightgray}
    \thead{APB Frequency} & \thead{\hyperref[fielddesc:EFUSETSUPA]{EFUSE\_TSUP\_A} \\ (> 6.669 ns)} & \thead{\hyperref[fielddesc:EFUSETPGM]{EFUSE\_TPGM} \\ (9-11 $\mu$s, usually 10 $\mu$s)} & \thead{\hyperref[fielddesc:EFUSETHPA]{EFUSE\_THP\_A} \\ (> 6.166 ns)} & \thead{\hyperref[fielddesc:EFUSETPGMINACTIVE]{EFUSE\_TPGM\_INACTIVE} \\ (> 35.96 ns)} \\
    \endhead
    \endfoot
    \endlastfoot
    \hline
    80 MHz & \hexprefix{}2 & \hexprefix{}320 & \hexprefix{}2 & \hexprefix{}4 \\
    \hline
    40 MHz & \hexprefix{}1 & \hexprefix{}190 & \hexprefix{}1 & \hexprefix{}2 \\
    \hline
    20 MHz & \hexprefix{}1 & \hexprefix{}C8 & \hexprefix{}1 & \hexprefix{}1 \\
    \hline
\end{longtable}

图 \ref{Figure:efuse_pgm} 中 A0 地址烧写 1，即 A0 对应的 eFuse 位为 1；A1 地址处于不烧写状态，A1 对应的 eFuse 位为 0。

\subsubsection{eFuse VDDQ 时序}\label{eFuse VDDQ Timing}
用户需要根据不同的 APB 时钟频率来配置
eFuse 的烧写电压 VDDQ 的时序参数，具体配置如下：

\begin{longtable}[c]{|c|c|c|c|}
    \caption{VDDQ 时序参数配置}
    \label{table:VDDQtiming}
    \\\hline
    \rowcolor{lightgray}
    \thead{APB Frequency} & \thead{\hyperref[fielddesc:EFUSEDACCLKDIV]{EFUSE\_DAC\_CLK\_DIV} \\ (> 1 $\mu$s)} & \thead{\hyperref[fielddesc:EFUSEPWRONNUM]{EFUSE\_PWR\_ON\_NUM} \\ (> EFUSE\_DAC\_CLK\_DIV*255)} & \thead{\hyperref[fielddesc:EFUSEPWROFFNUM]{EFUSE\_PWR\_OFF\_NUM}  \\ (> 3 $\mu$s)} \\
    \endfirsthead
    \hline
    \rowcolor{lightgray}
    \thead{APB Frequency} & \thead{\hyperref[fielddesc:EFUSEDACCLKDIV]{EFUSE\_DAC\_CLK\_DIV} \\ (> 1 $\mu$s)} & \thead{\hyperref[fielddesc:EFUSEPWRONNUM]{EFUSE\_PWR\_ON\_NUM} \\ (> EFUSE\_DAC\_CLK\_DIV*255)} & \thead{\hyperref[fielddesc:EFUSEPWROFFNUM]{EFUSE\_PWR\_OFF\_NUM}  \\ (> 3 $\mu$s)} \\
    \hline
    \endhead
    \endfoot
    \endlastfoot
    \hline
    80 MHz & \hexprefix{}A0 & \hexprefix{}A200 & \hexprefix{}100 \\
    \hline
    40 MHz & \hexprefix{}50 & \hexprefix{}5100 & \hexprefix{}80 \\
    \hline
    20 MHz & \hexprefix{}28 & \hexprefix{}2880 & \hexprefix{}40 \\
    \hline
\end{longtable}

\subsubsection{eFuse 读取时序}\label{eFuse-Read Timing}

图 \ref{Figure:efuse_read} 是 eFuse 读取的时序图。硬件提供 \hyperref[fielddesc:EFUSETSURA]{EFUSE\_TSUR\_A}、\hyperref[fielddesc:EFUSETRD]{EFUSE\_TRD} 和 \hyperref[fielddesc:EFUSETHRA]{EFUSE\_THR\_A} 三个寄存器用于控制 eFuse 的读时序。

\begin{figure}[H]
    \centering
    \fbox{\includegraphics[width=0.75\textwidth]{\modulefiles/figures/efuse_read_20191025}}
    \caption{eFuse 读取时序图}
    \label{Figure:efuse_read}
\end{figure}

用户需要根据不同的 APB 时钟频率来配置上述时序参数，具体请见表 \ref{table:Read_Timing_Configuration}。

\begin{longtable}[c]{ |c|c|c|c|}
    \caption{eFuse 读取时序参数配置}
    \label{table:Read_Timing_Configuration}
    \\\hline
    \rowcolor{lightgray}
    APB Frequency & \makecell[c]{\hyperref[fielddesc:EFUSETSURA]{EFUSE\_TSUR\_A} \\ (> 6.669 ns)} & \makecell[c]{\hyperref[fielddesc:EFUSETRD]{EFUSE\_TRD} \\ (> 35.96 ns)} & \makecell[c]{\hyperref[fielddesc:EFUSETHRA]{EFUSE\_THR\_A} \\ (> 6.166 ns)} \\
    \endfirsthead
    \hline
    \rowcolor{lightgray}
    APB Frequency & \makecell[c]{\hyperref[fielddesc:EFUSETSURA]{EFUSE\_TSUR\_A} \\ (> 6.669 ns)} & \makecell[c]{\hyperref[fielddesc:EFUSETRD]{EFUSE\_TRD} \\ (> 35.96 ns)} & \makecell[c]{\hyperref[fielddesc:EFUSETHRA]{EFUSE\_THR\_A} \\ (> 6.166 ns)} \\
    \endhead
    \endfoot
    \hline
    \endlastfoot
    \hline
    80 MHz & \hexprefix{}2 & \hexprefix{}4 & \hexprefix{}2 \\
    \hline
    40 MHz & \hexprefix{}1 & \hexprefix{}2 & \hexprefix{}1 \\
    \hline
    20 MHz & \hexprefix{}1 & \hexprefix{}1 & \hexprefix{}1 \\
\end{longtable}


\subsection{硬件模块使用参数}

硬件模块使用参数是通过电路连接实现的，用户无法干预这个过程。硬件使用的参数为表 \ref{table:BLOCK0_Parameter} 和 \ref{table:BLOCK1_10_Parameter}“硬件使用”一栏中标记为“Y”的参数。

\subsection{中断}
\begin{itemize}
\item 烧写完成中断：当 eFuse 烧写完成后，此中断被触发。如果要启动该中断信号，需将寄存器 \hyperref[fielddesc:EFUSEPGMDONEINTENA]{EFUSE\_PGM\_DONE\_INT\_ENA} 置 1。
\item 读取完成中断：当 eFuse 读取完成后，此中断被触发。如果要启动该中断信号，需将寄存器 \hyperref[fielddesc:EFUSEREADDONEINTENA]{EFUSE\_READ\_DONE\_INT\_ENA} 置 1。
\end{itemize}

\section{基地址}\label{efuse base address}
用户可以通过两个不同的寄存器基地址访问 eFuse 控制器，如表 \ref{table:efusebaseaddress} 所示。更多有关通过不同总线访问外设的信息，请参考章节 \ref{mod:sysmem} \textit{\nameref{mod:sysmem}}。
\newpage

\begin{longtable}[c]{ | l | C{3cm} | }
    \caption{eFuse 控制器基地址}
    \label{table:efusebaseaddress}
    \\\hline
    \rowcolor{lightgray}
    访问总线  & 基地址 \\\hline
    PeriBUS1 & \hexprefix{}3F41A000\\\hline
    PeriBUS2  &  \hexprefix{}6001A000\\\hline
\end{longtable}

\hypertarget{efuse-reg-summ}{}
\section{寄存器列表}

请注意，这里的地址是相对于 eFuse 控制器基地址的地址偏移量（相对地址）。请参阅章节 \ref{efuse base address} 获取有关 eFuse 控制器的基地址的信息。

\subfile{\modulefiles/efuse_summary_cn}

\section{寄存器}
\subfile{\modulefiles/efuse_reg_cn}
\end{document}
